\section{Introduction}\label{sec:introduction}

Let $\eta_\eps$ be a homogeneous Poisson process in $\R^2$ of intensity
$\eps^{-2}$. Color its Voronoi cells independently black or white, each
with probability $1/2$. For a bounded quadrilateral $Q$, a black crossing
means a connected black-cell path inside the closed quadrilateral joining
its two specified opposite sides. The cells are those of the full-plane
tessellation: the process is not resampled at the boundary, and no
boundary sites are wired together.

A nucleus is \emph{color-pivotal} for this event if changing only its cell's
color changes whether the crossing occurs. Write $N_\eps(Q)$ for the
number of such nuclei, including nuclei outside $Q$ when their cells
meet $Q$, and set $N_\eps=N_\eps([0,1]^2)$ for a left-to-right crossing.
All expectations below average both the Poisson geometry and the colors.

\begin{assumption}[Scalar Cardy limit]\label{ass:cardy}
For every bounded Jordan quadrilateral $Q$ with four distinct marked
boundary points, the critical crossing probability converges, as
$\eps\downarrow0$, to Cardy's conformally invariant value.
\end{assumption}

Here a Jordan quadrilateral is a Jordan domain with four marked boundary
points in cyclic order and a specified pair of opposite boundary arcs.
In the upper half-plane normalization with marked points $0,x,1,\infty$,
$0<x<1$, the crossing between $[0,x]$ and $[1,\infty]$ has Cardy value
\[
 \frac{\Gamma(2/3)}{\Gamma(1/3)\Gamma(4/3)}
 x^{1/3}\,{}_2F_1\!\left(\frac13,\frac23;\frac43;x\right).
\]
Assumption~\ref{ass:cardy} concerns one crossing event at a time. It
neither asserts a quantitative convergence rate nor assumes an
exploration or multi-crossing scaling limit.

\begin{theorem}\label{thm:main}
Under Assumption~\ref{ass:cardy}, there exists a constant
$c_V\in(0,\infty)$ such that
\[
 \lim_{\eps\downarrow0}\eps^{3/4}\E N_\eps=c_V.
\]
Consequently, the inverse expected pivotal count is asymptotic to
$c_V^{-1}\eps^{3/4}$.
\end{theorem}

The constant refers to the specified square, crossing event, intensity
convention, and color-pivotal count. The theorem is conditional on
Assumption~\ref{ass:cardy}; it does not prove Cardy's formula for the
Voronoi model.

\subsection{History and significance}
Cardy's formula~\cite{Cardy} predicts crossing probabilities at criticality
from the conformal shape of a planar domain. Smirnov proved the formula
and conformal invariance of crossing probabilities for critical site percolation on the triangular
lattice~\cite{Smirnov}. Schramm's stochastic Loewner evolution,
the exponent calculations of Lawler, Schramm, and Werner, and their
percolation application by Smirnov and Werner~\cite{Schramm,LSW,SW}
turned this conformal description into precise information about rare
connection events. In particular, the alternating four-arm exponent is $5/4$.
A bulk site is pivotal only when four alternating arms emerge from its
neighborhood, so this exponent predicts a pivotal count of order
$\eps^{-2}\eps^{5/4}=\eps^{-3/4}$.

The distinction between an exponent and an amplitude is essential here.
An estimate $\E N_\eps=\eps^{-3/4+o(1)}$ allows a slowly varying or
oscillating multiplicative correction. Sharp arm asymptotics, with a
limiting multiplicative constant, have been proved for critical
triangular-lattice percolation by Du, Gao, Li, and Zhuang~\cite{DGLZ}.
Theorem~\ref{thm:main} concerns a Poisson--Voronoi color-pivotal
observable, under a qualitative crossing-limit hypothesis, rather than
an unconditional triangular-lattice arm probability.

For Voronoi percolation, box-crossing and arm estimates~\cite{Tassion,Vanneuville} provide the
multiscale control needed to study critical and near-critical behavior;
we recall the precise estimates used here in Section~\ref{sec:inputs}.
Pivotal measures and their role in near-critical scaling limits were
developed for triangular-lattice percolation by Garban, Pete, and
Schramm~\cite{GPS}. The companion manuscript~\cite{V} derives a
quenched near-critical theory for Voronoi percolation from the scalar
Cardy hypothesis. In particular, it calibrates local pivotal
observables against the inverse expected pivotal count. Its
normalization remains defined through that count. Our task is to show
that this normalization is a pure power of the Euclidean mesh, with a
positive limiting coefficient.

Several ingredients have established predecessors. The passage from
local crossing data to joint crossing information uses the
quad-crossing framework and gluing theory of Schramm and
Smirnov~\cite{SS}. The annular exploration calculation uses Zhan's
annulus SLE~\cite{Zhan}. We state the precise inputs from~\cite{V}
before beginning the proof, and prove the additional annular and
signed-transfer estimates needed for the amplitude.

\subsection{The mechanism of the proof}
Work temporarily at unit Poisson intensity and enlarge the observation
window. Let
\[
 Q_0=[-1,1]\times[-1/2,1/2],\qquad
 f(R)=\Pp\{\text{a site inserted at $0$ is color-pivotal in $RQ_0$}\}.
\]
The inputs from~\cite{V} imply $f(R)=R^{-5/4+o(1)}$ and allow the final
answer to be recovered from the existence of
$\lim_{R\to\infty}R^{5/4}f(R)\in(0,\infty)$.

To improve the exponent to an amplitude, we compare \emph{differences}
of connection states across an annulus. A connection state records
which boundary color runs are joined through the interior, equivalently
the noncrossing pairing of the interfaces between those runs. The
transmitted difference vanishes unless at least four interfaces survive
across the annulus. With exactly four surviving interfaces, the two
possible interior states differ by a single scalar. This is the
source of a rank-one approximation for long-annulus transmission.

There is a second, stronger estimate. Adding one fair-colored Poisson
point produces a signed change that is invariant under color reversal
and rotations. In this symmetry class, the leading four-interface
contribution cancels. An explicit spectral calculation shows that the
remaining transmission decays faster than the inverse square of the
annulus scale ratio. This threshold is important: the point to be
inserted ranges over the two-dimensional plane, so the insertion
effect must be integrable in its location. An unsigned four-arm bound,
with exponent $5/4$, cannot supply that integrability.

The proof has three stages. First, Sections~\ref{sec:annuli}--\ref{sec:shape}
construct the annular transmission limits and establish their
rank-one asymptotic and signed rotational cancellation. Next,
Section~\ref{sec:transfer} makes these estimates uniform over arbitrary
signed boundary data and iterates them in the discrete Poisson model.
Finally, Section~\ref{sec:amplitude} differentiates $f$ with respect to
Poisson intensity. The integrable insertion estimate controls distant
insertions, while the rank-one estimate compares the normalized effect
of nearby insertions at two outer scales. Together they show that the
logarithmic derivative of $f$ has a summable error from $-5/4$.

The quantitative estimates do not come from a rate in
Assumption~\ref{ass:cardy}. We first choose a large but fixed annulus
ratio, use qualitative convergence to obtain the required contraction
beyond a sufficiently large starting scale, and only then iterate.
The uniformity over arbitrary signed boundary data, and the separation
between rank-one contraction and symmetry-enhanced decay, are the
technical features that may be useful for other amplitude problems.
